%==================================================================
\section{What the 2022 paper actually tested}

The last time this professor set this paper (22 September 2022; 60 marks, 2\,h\,30\,m, ``you may
answer any part of any question, but maximum you can score is 60''), every question was an
\textbf{algorithm-design} question. Here is the mapping onto this crash course:

\begin{center}
\begin{tabular}{@{}llll@{}}
\toprule
2022 Q & Topic & Technique & Here \\
\midrule
1(a) & Sort $A[i]\in\{1,2,3,4\}$ in $\Oh n$ & counting sort & \S6 facts \\
1(b) & In-place rearrange by a $0/1$ key & two-pointer partition & \S3 (partition step) \\
2(a) & Convex hull in $\Oh{n\log n}$ & sort $+$ scan & \S6 facts \\
2(b) & Convex hull is $\Om{n\log n}$ & reduction from sorting & \S6 \\
3(a) & $\E(\#\text{inversions})$ & linearity of expectation & Practice Q3 \\
3(b) & $T(n)=aT(n/b)+cn^k$, $a>b^k$ & unroll the geometric sum & \S2 \\
4 & Unit interval with most points & two-pointer sweep & Practice Q4 \\
5 & Subset sum closest to $x$ from below & DP & \S5 template \\
6 & Count pairs $A[i]>2A[j]$ in $\Oh{n\log n}$ & merge-sort variant & Practice Q5 \\
7 & $k$ numbers nearest the median in $\Oh n$ & median of medians & \S3 facts \\
\bottomrule
\end{tabular}
\end{center}

\begin{facts}{what this tells you about the exam}
\begin{itemize}
  \item \textbf{Nothing was pure bookwork.} Even the two theory parts (2(b), 3(b)) were
        \emph{proofs}, not statements.
  \item \textbf{Every algorithm question named a target complexity.} That target is the hint:
        $\Oh n$ means counting/two-pointer/selection; $\Oh{n\log n}$ means sort-then-scan or a
        merge-sort variant; $\Oh{n^2}$ means DP or all-pairs.
  \item \textbf{The 2022 paper contained no graph theory} --- but \S\S9--13 of your notes are
        entirely graph theory, so the syllabus has clearly widened. Prepare both. Practice
        Q7--Q9 below cover that half.
  \item \textbf{Always give all four parts:} algorithm, correctness, complexity, and the
        data structure that achieves it. The marks are split across them.
\end{itemize}
\end{facts}

%==================================================================
\section{Practice Set}

\noindent\emph{In the style and at the level of the 2022 paper. Total available: 105 marks; treat
any 60 of them as a full paper. Solutions begin in \S\ref{sec:sols} --- attempt first.}

\qhead{5+5}
\pt{a} $A$ is an unsorted array of size $n$ with $A[i]\in\{0,1,\dots,k-1\}$ for a fixed constant
$k$. Give an $\Oh n$-time, $\Oh1$-extra-space algorithm to sort $A$. What changes if
$k=\Th n$?

\pt{b} An array of $n$ items is given, each coloured Red, White or Blue. Rearrange the array
\emph{in place} in $\Oh n$ time and $\Oh 1$ extra space so that all Reds come first, then all
Whites, then all Blues. State the loop invariant.

\qhead{7+3}
\pt{a} Given $n$ points in the plane, design an $\Oh{n\log n}$ algorithm to compute the convex
hull. Prove it correct and justify the complexity.

\pt{b} Prove that any comparison-based algorithm computing the convex hull of $n$ points
requires $\Om{n\log n}$ time.

\qhead{5+5}
\pt{a} Let $A$ be a uniformly random permutation of $n$ distinct numbers. An \emph{inversion} is
a pair $(i,j)$ with $i<j$ and $A[i]>A[j]$; let $X$ be the number of inversions. Compute $\E(X)$.
Then give an $\Oh{n\log n}$ algorithm to count the inversions of a \emph{given} array.

\pt{b} Let $T(n)=a\,T(n/b)+c\,n^k$ with $T(1)=d$, $a,c>0$, $k,d\ge0$, $b\ge2$. Prove that if
$a<b^k$ then $T(n)=\Th{n^k}$.

\qhead{10}
$S=\langle a_1<a_2<\cdots<a_n\rangle$ is an increasing sequence of points on the real line, and
$k\le n$ is a positive integer. Give an $\Oh n$-time algorithm that finds the \textbf{shortest}
interval containing at least $k$ points of $S$. Prove correctness.

\qhead{10}
Let $A$ be an array of $n$ distinct reals. Call $(i,j)$ a \emph{heavy pair} if $i<j$ and
$A[i]>3A[j]+5$. Design an $\Oh{n\log n}$ algorithm to count the heavy pairs.

\qhead{10}
There are $n$ jobs; job $i$ occupies the interval $[s_i,f_i]$ and has weight $w_i>0$. Find a
maximum-weight set of pairwise non-overlapping jobs in $\Oh{n\log n}$. Prove optimal substructure.
What does this problem become in the language of \S12, and why is that significant?

\qhead{10}
Let $T$ be a tree on $n$ vertices. Give an $\Oh n$ algorithm that computes the eccentricity
$\mathrm{ecc}(v)=\max_u d(v,u)$ for \emph{every} vertex $v$ simultaneously. Prove it correct.

\qhead{5+5}
\pt{a} Prove that a simple planar graph with $n\ge3$ vertices has at most $3n-6$ edges, and
deduce that $K_5$ is not planar.

\pt{b} Prove that a \emph{triangle-free} simple planar graph with $n\ge3$ vertices has at most
$2n-4$ edges, and deduce that $K_{3,3}$ is not planar.

\qhead{10}
$G$ is a chordal graph on $n$ vertices and $m$ edges, and a perfect elimination ordering of $G$
is given. Give an $\Oh{n+m}$ algorithm that computes $\chrom(G)$ and $\clq(G)$, and prove that
they are equal.

\qhead{10}
Given a set $S$ of $n$ distinct reals, decide whether there exist four \emph{distinct} elements
$x,y,z,w\in S$ with $x+y=z+w$. Give an algorithm and analyse it. (Aim for $\Oh{n^2}$ up to the
cost of a dictionary.)

\qhead{10}
You are given the \textbf{preorder} traversal of a binary search tree on $n$ distinct keys.
Reconstruct the tree in $\Oh n$ time and prove correctness. Explain why the analogous problem for
a general binary tree is impossible.

\qhead{5}
Prove that any comparison-based algorithm that merges two sorted arrays of size $n$ each requires
at least $2n-1$ comparisons in the worst case.

%==================================================================
\section{Solutions to the Practice Set}\label{sec:sols}

\setcounter{pq}{0}

\qhead{5+5}
\pt{a} \textbf{Counting sort.} Keep an array $C[0..k-1]$ of counters (size $k$, a constant, hence
$\Oh1$ space). One pass over $A$ incrementing $C[A[i]]$; then one pass writing $C[0]$ copies of
$0$, then $C[1]$ copies of $1$, and so on. Time $\Oh{n+k}=\Oh n$; extra space $\Oh k=\Oh1$.

\emph{Correctness:} the output is non-decreasing by construction, and it is a permutation of the
input because each value $v$ is written exactly $C[v]=\#\{i:A[i]=v\}$ times.

\emph{No contradiction with the $\Om{n\log n}$ bound:} this is not a comparison sort --- it uses
the values as array indices.

\emph{If $k=\Th n$:} the time is still $\Oh{n+k}=\Oh n$, but the counter array now needs $\Th n$
space, so the $\Oh1$-space claim fails. It becomes an $\Oh n$-time, $\Oh n$-space algorithm.

\pt{b} \textbf{Dutch national flag, three pointers.} Maintain $\mathit{lo}=0$, $\mathit{mid}=0$,
$\mathit{hi}=n-1$ and loop while $\mathit{mid}\le \mathit{hi}$:
\begin{itemize}
  \item $A[\mathit{mid}]$ Red: swap $A[\mathit{lo}]\leftrightarrow A[\mathit{mid}]$; $\mathit{lo}{+}{+}$; $\mathit{mid}{+}{+}$.
  \item $A[\mathit{mid}]$ White: $\mathit{mid}{+}{+}$.
  \item $A[\mathit{mid}]$ Blue: swap $A[\mathit{mid}]\leftrightarrow A[\mathit{hi}]$; $\mathit{hi}{-}{-}$ (do \emph{not} advance $\mathit{mid}$).
\end{itemize}
\textbf{Loop invariant:} $A[0,\mathit{lo})$ all Red, $A[\mathit{lo},\mathit{mid})$ all White,
$A(\mathit{hi},n)$ all Blue, and $A[\mathit{mid},\mathit{hi}]$ unexamined.

\emph{Correctness:} each branch restores the invariant --- the Blue case cannot advance
$\mathit{mid}$ because the element swapped in from $\mathit{hi}$ is still unexamined. On
termination $\mathit{mid}>\mathit{hi}$, so the unexamined region is empty and the three blocks
tile the array.

\emph{Complexity:} every iteration either increments $\mathit{mid}$ or decrements $\mathit{hi}$,
so the quantity $\mathit{hi}-\mathit{mid}$ strictly decreases: at most $n$ iterations, each
$\Oh1$. Time $\Oh n$, extra space $\Oh1$ (three indices). The same three-way partition is exactly
what a robust quicksort/quickselect uses when duplicates are present.

\qhead{7+3}
\pt{a} \textbf{Andrew's monotone chain.}
\begin{enumerate}
  \item Sort the points lexicographically by $(x,y)$: $\Oh{n\log n}$.
  \item \textbf{Lower hull:} scan left to right maintaining a stack; before pushing $p$, while the
        top two stack points $q,r$ and $p$ do \emph{not} make a counter-clockwise turn (i.e.\
        the cross product $(r-q)\times(p-q)\le0$), pop.
  \item \textbf{Upper hull:} the same scan on the reversed order.
  \item Concatenate, dropping the duplicated endpoints.
\end{enumerate}
\emph{Correctness:} a point is a vertex of the lower hull iff it makes a strict left turn with its
hull-neighbours. The stack maintains the invariant ``the current stack is the lower hull of the
points scanned so far'': popping removes exactly the points that the new point makes non-convex,
and no popped point can return to the hull, since it lies below the segment joining two points
that remain. Doing this in both directions gives the full hull.

\emph{Complexity:} $\Oh{n\log n}$ for the sort; the scan pushes each point once and pops it at
most once, so it is $\Oh n$ amortised. Total $\Oh{n\log n}$, which by (b) is optimal.

\pt{b} \textbf{Reduction from sorting.} Given reals $x_1,\dots,x_n$, build the points
$(x_i,x_i^2)$ in $\Oh n$ time. All of them lie on the parabola $y=x^2$, which is strictly convex,
so \emph{every} input point is a hull vertex. A hull algorithm outputs the vertices in cyclic
order; starting at the leftmost and reading along the lower chain gives $x_1,\dots,x_n$ in
increasing order --- an $\Oh n$ extraction.

So a hull algorithm running in $T(n)$ would sort in $T(n)+\Oh n$. Since comparison sorting is
$\Om{n\log n}$ (\S6), $T(n)+\Oh n=\Om{n\log n}$, hence $T(n)=\Om{n\log n}$. \qed

\qhead{5+5}
\pt{a} \textbf{Expectation.} For each pair $i<j$ let $X_{ij}=\mathbf 1\{A[i]>A[j]\}$, so
$X=\sum_{i<j}X_{ij}$. In a uniformly random permutation the two orderings of any fixed pair are
equally likely, so $\E(X_{ij})=\tfrac12$. By linearity (no independence needed):
\[
  \E(X)=\binom n2\cdot\frac12=\boxed{\frac{n(n-1)}{4}} .
\]

\textbf{Counting algorithm --- modified merge sort.} Sort recursively, and during the merge of
the sorted halves $L$ and $R$: whenever an element of $R$ is output while $t$ elements of $L$
remain unconsumed, add $t$ to the counter (that element of $R$ is inverted with exactly those $t$
elements).

\emph{Correctness:} every pair $(i,j)$ with $i<j$ has both indices in $L$, both in $R$, or split.
The first two are counted by the recursive calls, the third exactly once during this merge, and
$L$ holds precisely the smaller indices. So each inverted pair is counted once.

\emph{Complexity:} $T(n)=2T(n/2)+\Oh n=\Oh{n\log n}$ (Master Theorem, Case 2).

\pt{b} Unroll for $n=b^L$, $L=\log_b n$, exactly as in \S2:
\[
  T(n)=c\,n^k\sum_{i=0}^{L-1}\left(\frac{a}{b^k}\right)^{\!i}+d\,n^{\log_b a}.
\]
Put $r=a/b^k$. Now $a<b^k$ gives $r<1$, so the geometric series \emph{converges}:
\[
  \sum_{i=0}^{L-1}r^i\ \le\ \sum_{i=0}^{\infty}r^i=\frac{1}{1-r}=\Th1 ,
\]
and it is at least $1$. Hence the first term is $\Th{n^k}$. For the second, $a<b^k$ gives
$\log_b a<k$, so $n^{\log_b a}=o(n^k)$. Therefore
\[
  T(n)=\Th{n^k}+o(n^k)=\boxed{\Th{n^k}} .
\]
(The lower bound is immediate anyway: the root alone contributes $c\,n^k$.) \qed
\emph{Interpretation:} $r<1$ means the work \emph{shrinks} geometrically down the tree, so the
root dominates --- Master Theorem Case 3.

\qhead{10}
\textbf{Algorithm.} For $i=1,\dots,n-k+1$ compute $a_{i+k-1}-a_i$ and return the minimum (and the
interval attaining it). $\Oh n$ time, $\Oh 1$ space.

\emph{Correctness.} Let $I$ be any interval containing at least $k$ points of $S$. Since $S$ is
sorted, the points of $S$ inside $I$ are \textbf{consecutive}: $a_i,a_{i+1},\dots,a_j$ with
$j-i+1\ge k$. Then
\[
  |I|\ \ge\ a_j-a_i\ \ge\ a_{i+k-1}-a_i,
\]
the second inequality because $j\ge i+k-1$ and the $a$'s increase. So no interval beats the best
window. Conversely each window $[a_i,a_{i+k-1}]$ genuinely contains the $k$ points
$a_i,\dots,a_{i+k-1}$, so the minimum over windows is attained. \qed

\emph{Remark.} The 2022 paper asked the dual --- \emph{fixed} length $1$, maximise the count. That
one needs a two-pointer sweep: advance $j$ while $a_j\le a_i+1$, take $j-i$, then advance $i$.
Both pointers only move forward, so it is $\Oh n$. Know both directions.

\qhead{10}
\textbf{Algorithm --- merge-sort variant.} \code{CountAndSort}$(A[\ell..r])$:
\begin{enumerate}
  \item If $\ell\ge r$, return $0$.
  \item Split at the midpoint; recursively count-and-sort each half, obtaining sorted $L$, $R$
        and their counts.
  \item \textbf{Cross count.} Both halves are now sorted ascending. Walk $j$ over $R$ in
        increasing order; maintain a pointer $p$ into $L$ marking the first position with
        $L[p]>3R[j]+5$. Add $|L|-p$ to the count. Since $3R[j]+5$ increases with $j$, $p$ is
        non-decreasing --- one monotone sweep.
  \item Merge $L$ and $R$ normally and return the total count.
\end{enumerate}
\emph{Correctness.} Every pair $(i,j)$ with $i<j$ lies wholly in the left half, wholly in the
right half, or is split with $i$ in the left and $j$ in the right --- because merge sort splits by
\emph{index}, so all left-half indices precede all right-half indices. The first two cases are
handled by recursion; the third is counted in step 3, which for each $j$ counts exactly the
$L$-elements exceeding $3A[j]+5$. Sorting the halves is harmless: the count for a split pair
depends only on the two \emph{values}, not on their order within a half.

\emph{Complexity.} Steps 3 and 4 are each a single linear sweep, so
$T(n)=2T(n/2)+\Oh n=\boxed{\Oh{n\log n}}$ by the Master Theorem, Case 2.

\emph{Remark.} The same skeleton solves 2022 Q6 ($A[i]>2A[j]$) and inversion counting
($A[i]>A[j]$): only the threshold function changes. It must be \textbf{monotone} in $A[j]$ for the
single-pointer sweep to be valid.

\qhead{10}
\textbf{Weighted interval scheduling.}
\begin{enumerate}
  \item Sort jobs by finishing time so $f_1\le f_2\le\cdots\le f_n$: $\Oh{n\log n}$.
  \item For each $j$ compute $p(j)=\max\{i<j: f_i\le s_j\}$ (or $0$), by binary search in the
        sorted finish times: $\Oh{\log n}$ each, $\Oh{n\log n}$ total.
  \item DP with $\mathrm{opt}(0)=0$ and
        \[
          \mathrm{opt}(j)=\max\Bigl\{\,w_j+\mathrm{opt}(p(j)),\ \ \mathrm{opt}(j-1)\,\Bigr\}.
        \]
  \item Answer $\mathrm{opt}(n)$; recover the set by tracing which branch achieved each maximum.
\end{enumerate}
\emph{Optimal substructure.} Consider an optimal solution $O_j$ for jobs $1,\dots,j$.
\begin{itemize}
  \item If job $j\in O_j$: no job $i$ with $p(j)<i<j$ can be in $O_j$, since every such job
        finishes after $s_j$ and starts before $f_j$, hence overlaps $j$. So
        $O_j\setminus\{j\}$ is a feasible solution on $1,\dots,p(j)$, and it must be optimal
        there --- otherwise replacing it by a better one and re-adding $j$ would beat $O_j$.
        Value $w_j+\mathrm{opt}(p(j))$.
  \item If job $j\notin O_j$: $O_j$ is a feasible solution on $1,\dots,j-1$ and must be optimal
        there by the same exchange argument. Value $\mathrm{opt}(j-1)$.
\end{itemize}
These two cases are exhaustive, so the maximum of the two values is $\mathrm{opt}(j)$; induction
on $j$ completes the proof. \qed

\emph{Complexity.} $\Oh{n\log n}$ (sort $+$ binary searches) $+\ \Oh n$ (DP) $=\boxed{\Oh{n\log n}}$.

\emph{The connection.} Build the interval graph $G$ whose vertices are the jobs, with an edge
between overlapping jobs. A set of pairwise non-overlapping jobs is exactly an
\textbf{independent set} of $G$, so this is \textbf{maximum-weight independent set on an interval
graph}. That is significant because MWIS is NP-hard on general graphs, yet interval graphs are
chordal, hence perfect, and the structure (here: a linear order by finishing time) makes the
problem solvable in $\Oh{n\log n}$ --- the same phenomenon as \S13, where a PEO makes maximum
independent set easy on chordal graphs.

\qhead{10}
\textbf{Algorithm.} Three BFS runs.
\begin{enumerate}
  \item BFS from an arbitrary $s$; let $u$ be a farthest vertex.
  \item BFS from $u$, recording $d_u(\cdot)$; let $v$ be a farthest vertex. Then $(u,v)$ is a
        diametral pair.
  \item BFS from $v$, recording $d_v(\cdot)$.
  \item Output $\mathrm{ecc}(x)=\max\{d_u(x),d_v(x)\}$ for every $x$.
\end{enumerate}
\emph{Correctness.} Step 2 gives a diametral pair by the two-BFS trick (\S9). It remains to show:
for \emph{every} $x$ and any diametral pair $(u,v)$,
$\mathrm{ecc}(x)=\max\{d(x,u),d(x,v)\}$.

``$\ge$'' is trivial. For ``$\le$'', let $P$ be the (unique) $u$--$v$ path, let $w$ be any vertex,
and let $m$, $m_w$ be the vertices of $P$ closest to $x$ and to $w$ respectively.
First, since $d(u,w)=d(u,m_w)+d(m_w,w)\le\mathrm{diam}=d(u,m_w)+d(m_w,v)$, we get
$d(m_w,w)\le d(m_w,v)$, and symmetrically $d(m_w,w)\le d(m_w,u)$.
\begin{itemize}
  \item If $m=m_w$: $d(x,w)=d(x,m)+d(m,w)\le d(x,m)+\max\{d(m,u),d(m,v)\}=\max\{d(x,u),d(x,v)\}$.
  \item If $m\ne m_w$, say $m$ lies between $u$ and $m_w$ on $P$. Then
        \[
          d(x,v)=d(x,m)+d(m,v)=d(x,m)+d(m,m_w)+d(m_w,v)\ \ge\ d(x,m)+d(m,m_w)+d(m_w,w)=d(x,w).
        \]
        (The symmetric case gives $d(x,u)\ge d(x,w)$.)
\end{itemize}
Either way $d(x,w)\le\max\{d(x,u),d(x,v)\}$ for every $w$, proving the claim. \qed

\emph{Complexity.} Three BFS runs on a tree: $\boxed{\Oh{n}}$ (since $|E|=n-1$).

\qhead{5+5}
\pt{a} Every face of a plane drawing of a simple graph with $n\ge3$ is bounded by at least $3$
edges, and every edge borders exactly $2$ faces. Counting edge--face incidences two ways,
$2E\ge3F$, i.e.\ $F\le\tfrac23E$. Substituting into Euler's formula $n-E+F=2$:
\[
  2=n-E+F\ \le\ n-E+\tfrac23E=n-\tfrac13E
  \quad\Longrightarrow\quad E\le3n-6 .
\]
\textbf{$K_5$:} $n=5$, $E=\binom52=10$, but $3n-6=9<10$. So $K_5$ violates the bound and is not
planar. \qed

\pt{b} If $G$ is triangle-free then every face is bounded by at least $4$ edges, so $2E\ge4F$,
i.e.\ $F\le\tfrac12E$. Euler again:
\[
  2=n-E+F\le n-E+\tfrac12E=n-\tfrac12E
  \quad\Longrightarrow\quad E\le2n-4 .
\]
\textbf{$K_{3,3}$:} it is bipartite, hence triangle-free, with $n=6$ and $E=9$; but
$2n-4=8<9$. So $K_{3,3}$ is not planar. \qed

\emph{Note:} the first bound alone does \emph{not} rule out $K_{3,3}$ --- $3n-6=12\ge9$. That is
exactly why the triangle-free refinement is needed, and it is a favourite exam trap.

\qhead{10}
\textbf{Algorithm.} Let $v_1,\dots,v_n$ be the given PEO. Process it in order; for each $i$
compute
\[
  d_i=\bigl|N(v_i)\cap\{v_1,\dots,v_{i-1}\}\bigr|
\]
and give $v_i$ the lowest-indexed colour not used by those earlier neighbours. Output
$\max_i(d_i+1)$ as both $\chrom(G)$ and $\clq(G)$, together with the colouring.

\emph{Complexity.} Computing all the $d_i$ is one scan of the adjacency lists with a position
array, $\Oh{n+m}$. Choosing $v_i$'s colour needs a scratch array of size $d_i+1$ marking the
colours of its earlier neighbours, so the total colouring work is $\sum_i\Oh{d_i+1}=\Oh{n+m}$.
Total $\boxed{\Oh{n+m}}$.

\emph{Proof that $\chrom=\clq$.}
\begin{itemize}
  \item \textbf{Upper bound on $\chrom$.} When $v_i$ is coloured, its earlier neighbours form a
        \emph{clique} (PEO property) of size $d_i$, so they carry $d_i$ \emph{distinct} colours,
        and some colour in $\{1,\dots,d_i+1\}$ is free. Hence the greedy uses at most
        $1+\max_i d_i$ colours, so $\chrom(G)\le1+\max_i d_i$.
  \item \textbf{Lower bound on $\clq$.} For the index $i$ attaining the maximum, the set
        $\{v_i\}\cup\bigl(N(v_i)\cap\{v_1,\dots,v_{i-1}\}\bigr)$ is a clique of size $d_i+1$
        (the earlier neighbours are mutually adjacent, and $v_i$ is adjacent to all of them). So
        $\clq(G)\ge1+\max_i d_i$.
  \item \textbf{Universal inequality.} $\chrom(G)\ge\clq(G)$ always.
\end{itemize}
Chaining: $\clq\le\chrom\le1+\max_i d_i\le\clq$, so all are equal:
\[
  \boxed{\chrom(G)=\clq(G)=1+\max_i d_i .}
\]
\qed \emph{This is the constructive proof that chordal graphs are perfect} (the argument applies
verbatim to every induced subgraph, since an induced subgraph of a chordal graph is chordal and
inherits a PEO).

\qhead{10}
\textbf{Key observation.} If $\{x,y\}$ and $\{z,w\}$ are \emph{distinct} pairs with $x+y=z+w$,
they are automatically \textbf{disjoint}: if they shared an element, cancelling it would force
the other two to be equal, impossible among distinct elements. So we need only find two distinct
pairs with equal sums --- distinctness of all four is then free.

\textbf{Algorithm.} Form all $\binom n2$ pairwise sums, each tagged with its pair. Insert them
into a dictionary keyed by the sum. If any key is hit twice, output the two pairs; if no
collision occurs, answer ``no''.

\emph{Correctness.} A ``yes'' instance is precisely a repeated sum, by the key observation. The
algorithm reports exactly that.

\emph{Complexity.} $\binom n2=\Th{n^2}$ sums to generate.
\begin{itemize}
  \item With a hash table: $\Oh{n^2}$ expected time, $\Oh{n^2}$ space.
  \item Deterministically, by sorting the sums and scanning for adjacent duplicates:
        $\Oh{n^2\log n}$ time (which is $\Oh{n^2\log n}=\Oh{n^2\cdot\log n}$; note
        $\log\binom n2=\Th{\log n}$), $\Oh{n^2}$ space.
\end{itemize}
\emph{Remark.} One can stop as soon as a collision appears, so on ``yes'' instances the algorithm
often terminates far earlier --- but the worst case (a Sidon set, where all pairwise sums are
distinct) forces all $\Th{n^2}$ sums to be examined.

\qhead{10}
\textbf{Why preorder alone suffices for a BST.} The inorder traversal of a BST with distinct keys
is forced --- it is the keys in increasing order. So preorder $+$ inorder is available, and by the
standard reconstruction (\S4) that determines the tree uniquely. Sorting would cost
$\Oh{n\log n}$; the bounds method below achieves $\Oh n$.

\textbf{Algorithm (bounds method).} Keep a global index $\mathit{idx}=0$ into the preorder array
$P$. Define
\[
  \code{Build}(\mathit{lo},\mathit{hi}):
\]
\begin{enumerate}
  \item If $\mathit{idx}=n$, or $P[\mathit{idx}]\notin(\mathit{lo},\mathit{hi})$, return
        \code{null}.
  \item Let $\kappa=P[\mathit{idx}]$; create a node with key $\kappa$; $\mathit{idx}{+}{+}$.
  \item Left child $\leftarrow\code{Build}(\mathit{lo},\kappa)$;
        right child $\leftarrow\code{Build}(\kappa,\mathit{hi})$.
  \item Return the node.
\end{enumerate}
Call \code{Build}$(-\infty,+\infty)$.

\emph{Correctness.} By the BST property, in the preorder sequence the root $\kappa$ is followed
first by the entire preorder of its left subtree --- all of whose keys are $<\kappa$ --- and then
by the preorder of its right subtree, all of whose keys are $>\kappa$. So the split point is
exactly the first position after $\kappa$ carrying a key $>\kappa$, which is precisely what the
bound test $P[\mathit{idx}]\in(\mathit{lo},\mathit{hi})$ detects. Induction on subtree size gives
correctness.

\emph{Complexity.} $\mathit{idx}$ advances exactly once per node, and each node generates $\Oh1$
calls, so the total number of calls (including those returning \code{null}) is $\Oh n$, each
doing $\Oh1$ work: $\boxed{\Oh n}$.

\emph{Why a general binary tree is impossible.} With $n$ distinct keys there are far more binary
trees than preorder sequences can distinguish: concretely, a preorder ``$a$, then $b$'' is
consistent with $b$ being either the left child or the right child of $a$, two genuinely
different trees producing identical preorder output. Hence no algorithm can reconstruct. The BST
case escapes only because the extra ordering constraint supplies the missing inorder information.

\qhead{5}
\textbf{Adversary argument.} Consider inputs in which the correct merged output alternates
between the two arrays:
\[
  a_1<b_1<a_2<b_2<\cdots<a_n<b_n .
\]
The merged sequence has $2n$ elements, hence $2n-1$ adjacent pairs, and each adjacent pair
consists of one element from each array (by the alternation).

Suppose some algorithm halts having never compared some adjacent pair $(u,v)$ from different
arrays, where $u$ immediately precedes $v$ in the true output. Build a second input identical
except that the values of $u$ and $v$ are swapped. This second input is still a legal pair of
sorted arrays (swapping two adjacent values across the two arrays preserves the sorted order
within each array), and it is consistent with \emph{every} comparison the algorithm actually made
--- the only comparison whose outcome would change is the one between $u$ and $v$, which was never
performed. Yet the correct output differs, since $u$ and $v$ are transposed. So the algorithm
gives a wrong answer on one of the two inputs.

Hence a correct algorithm must compare all $2n-1$ adjacent pairs, requiring at least $2n-1$
comparisons in the worst case. \qed

\emph{Matching upper bound:} standard merging uses at most $2n-1$ comparisons, so the bound is
tight.

\vfill
\begin{center}\small\itshape
Sections 1--13 follow the lecture notes exactly; the practice set is modelled on the
mid-semester examination of 22 September 2022.
\end{center}
